%%% beispiel.tex — zeigt jedes Element der Klasse genau einmal.
%%%
%%%   cd beispiel && xelatex beispiel.tex     (dreimal)
%%%
%%% Vorher muessen die Grafiken und Schriften da sein:
%%%   python3 ../werkzeuge/aufbereiten.py "/pfad/zu/Scriptorium Aventuris v4"
%%%   python3 ../werkzeuge/pruefen.py
%%%
%%% Dieses Dokument ist zugleich der Probelauf. Die Abschnitte sind einzeln
%%% auskommentierbar, um einen Fehler einzukreisen. Die Liste der Stellen, die
%%% beim ersten Lauf zu pruefen sind, steht in doku/ELEMENTE.md.

\documentclass[raster,rasterzeigen]{dsa5latex}

% Die Klasse sucht Grafiken und Schriften relativ zum Arbeitsverzeichnis.
% Von beispiel/ aus liegen sie eine Ebene hoeher.
\graphicspath{{../}}

% Der Abenteuertitel steht im Kolumnentitel neben der Seitenzahl.
\dsaAbenteuertitel{Beispielabenteuer}

\begin{document}

%%% ==================================================================
%%% Umschlag
%%% ==================================================================
% Braucht ein Coverbild. Ohne eigenes: die Aventurienkarte nehmen, die das
% Aufbereitungswerkzeug mitliefert.
\dsaUmschlagVorne{grafiken/aventurienkarte}{%
  \dsaTitelZeile{Ein Beispiel}%
  \dsaTitelZeile{in zwei Zeilen}}

%%% ==================================================================
%%% Impressum
%%% ==================================================================
\begin{dsaImpressumseite}

% Rubriken nach dem Vorbild einer gesetzten Scriptorium-Veroeffentlichung.
\dsaImpressumsblock{Autor}{[Name]}
\dsaImpressumsblock{Redaktion}{[Name]}
\dsaImpressumsblock{Lektorat}{[Name]}
\dsaImpressumsblock{Korrektorat}{[Name]}
\dsaImpressumsblock{Künstlerische Leitung}{[Name]}
\dsaImpressumsblock{Coverbild}{[Name]}
\dsaImpressumsblock{Satz, Layout und Gestaltung}{[Name]}
\dsaImpressumsblock{Innenillustrationen und Pläne}{[Namen]}
\dsaImpressumsblock{Version}{1.0}

\dsaRechtevermerk{2026}{[Name oder Firma]}

\end{dsaImpressumseite}

\dsaInhalt

%%% ==================================================================
\dsakapitel{Satzspiegel und Raster}
%%% ==================================================================

Diese Seite ist mit \verb|rasterzeigen| gebaut. Die feinen Linien sind das
Grundlinienraster, 12 pt Abstand, erste Linie 12,7 mm unter der
Papieroberkante. Die staerkere Linie ist die Oberkante des Satzspiegels bei
24 mm.

\textbf{Was zu pruefen ist:} Jede Textzeile muss auf einer feinen Linie
sitzen, und die Zeilen beider Spalten muessen auf derselben Hoehe liegen.

Sollmasse: Satzbreite 166 mm, Spaltenbreite 80,5 mm, Spaltensteg 5 mm, 59
Zeilen je Spalte. Zugleich der Test fuer die deutsche Silbentrennung:
Grundlinienraster, Abenteuerzusammenfassung, Meisterinformationen,
Aufzaehlungszeichen, Seitenhintergrundgrafik, Spaltenzwischenraum.

\dsaunterkapitel{Unterkapitel, zentriert}
\dsaabschnitt{Abschnitt, Gentium Basic fett 13 pt}
\dsaunterabschnitt{Unterabschnitt, Gentium Basic fett 10 pt}

Alle vier Groessen sind im Klartext des Verlags belegt.

\dsaAbschnittGelb{Abschnitt mit Pergamentverlauf}

\clearpage

%%% ==================================================================
\dsakapitel{Fliesstext und Marken}
%%% ==================================================================

\dsaEinfuehrung{Einfuehrungstext, kursiv mit Laufweite plus zehn.}

\dsaStimmung{Wer zaehlt, der irrt sich nur einmal. Wer nicht zaehlt, irrt
sich immer.}{Ein Sprichwort, aufgeschrieben irgendwo}

\dsaZitat{Freistehendes Zitat, Grauwert 404040, kursiv, zentriert.}

Gewoehnlicher Fliesstext im Blocksatz, Gentium Basic 10 pt auf 12 pt. Ein
Verweis steht hochgestellt hinter dem Regelelement, etwa
Trampeln\dsaBand{ABE}{8}. Zahlenangaben stehen im Textmodus:
\dsaFormel{1W6+4} Schaden, ein Modifikator von \dsaMod{-}{2}, eine Tragkraft
von KK\dsaMal 2. \dsaKapitaelchen{Nachgebildete Kapitaelchen} — keine der
Schriften hat echte.

Meisterinformation beginnt hier \dsaAugeSchwarz{} und endet im Kasten
\dsaAugeWeiss{}. Fokusregel: \dsaFokusregel{Neu}. Gegnerabstufung:
\dsaBauer{} \dsaSpringer{} \dsaTurm{} \dsaKoenig{}. Rautenskalen, in TikZ
gezeichnet: \dsaRautenRot{2}{18mm} und \dsaRautenGruen{4}{18mm}.

\dsaFiole{l}
Die Fiole am linken Rand markiert einen Vorschlag, die Szene leichter zu
machen. Sie ist 9,1 mm gross, wie im Baukasten.

\dsaSchaedel{r}
Der Totenkopf markiert das Gegenteil. Ebenfalls 9,1 mm.

\begin{dsaliste}
\item Aufzaehlung mit der DSA-Raute
\item zweite Zeile
\end{dsaliste}

\begin{dsalisteblau}
\item Blaue Raute
\end{dsalisteblau}

\begin{dsaWerteabsatz}
Werteabsatz mit haengendem Einzug von 8,504 pt. Die zweite Zeile ruecken ein.
\end{dsaWerteabsatz}

\dsaVorlesetext{Vorlesetext mit der Zierleiste darueber.}

\clearpage

%%% ==================================================================
\dsakapitel{Kaesten}
%%% ==================================================================

\begin{dsaPergamentKlein}
Pergament klein, 85,5 mal 50,7 mm, 12 Rastereinheiten. Ragt 2,5 mm je Seite
ueber die Spalte — so gewollt, der Zierrand liegt ausserhalb des
Textbereichs.
\end{dsaPergamentKlein}

\begin{dsaWerteMittel}
Wertekasten mittel, 86,9 mal 103,8 mm.

\begin{dsaWerteblock}
MU 11 \quad KL 13 \quad IN 14 \quad CH 10\\
LeP 30 \quad INI \dsaFormel{12+1W6}
\end{dsaWerteblock}
\end{dsaWerteMittel}

\clearpage

% Das Portraitbild kommt als Schluessel, nicht als Befehl im Inhalt.
\begin{dsaWerteMittelPortrait}[dsaportrait=grafiken/aventurienkarte]
Wertekasten mittel mit Portrait, 94,0 mal 104,0 mm. Das Medaillon tritt
13,5 mm zur Aussenseite heraus. Vorlaeufig nur die rechte Lage.
\end{dsaWerteMittelPortrait}

% Der Gegner- und Kreaturkasten: dieselbe Vorlagengrafik, aber in freier
% Hoehe aus drei Scheiben gebaut.
\begin{dsaWerteFreiPortrait}[dsaportrait=grafiken/aventurienkarte]{18}
{\bfseries Wolfsratte}

\begin{dsaWerteblock}
MU 11 \quad KL 2 \quad IN 13 \quad CH 8\\
FF 10 \quad GE 14 \quad KO 12 \quad KK 11
\end{dsaWerteblock}
\dsaFeld{LeP}{18}
\dsaFeld{INI}{\dsaFormel{13+1W6}}
\dsaFeld{Biss}{AT 12 TP \dsaFormel{1W6+1} RW kurz}
\dsaFeld{RS/BE}{1/0}
\dsaFeld{Groesse}{klein}
\dsaFeld{Typus}{Tier}
\end{dsaWerteFreiPortrait}

\begin{dsaMeisterMaske}
Meisterkasten mit angesetzter Maske, 84,0 mal 108,0 mm, Text weiss. Der
Regelfall im Text, weil er fast Spaltenbreite hat.
\end{dsaMeisterMaske}

\begin{dsaMeisterSchmal}
Meisterkasten, 59,9 mal 114,2 mm — hochkant.
\end{dsaMeisterSchmal}

\clearpage
\onecolumn

\begin{dsaMeisterBreit}
Maskenfeld breit, 177,9 mal 114,2 mm, breiter als die Satzbreite. Einspaltig.
\end{dsaMeisterBreit}

\twocolumn

\begin{dsaKastenFrei}{14}
Kasten in freier Hoehe, 14 Rastereinheiten. Pergamentflaeche in der
gewuenschten Hoehe, Zierleisten oben und unten in wahrer Groesse — die
Bauweise des Verlags.
\end{dsaKastenFrei}

\dsaMeistermaske[2]{Die fertige Meistermaske aus dem Baukasten, mittlere
Hoehe. Sie enthaelt die gestrichelte Verbindung schon.}

\clearpage

%%% ==================================================================
\dsakapitel{Bilder, Umfluss, Tabellen}
%%% ==================================================================

\dsaBildSpalte{grafiken/aventurienkarte}{10}

Darunter laeuft der Text weiter und sitzt wieder auf der Grundlinie, weil das
Bild eine ganze Zahl Rastereinheiten belegt.

\dsaabschnitt{Kreis und Freiform}

\dsaBildKreis{25mm}{grafiken/aventurienkarte}
\dsaBildForm{30mm}{20mm}{(0,0) (1,0.6) (1.6,-0.3) (0.6,-1)}{grafiken/aventurienkarte}

\dsaabschnitt{Umfluss mit festem Einzug}

\dsaUmfluss{l}{2}{4}{25mm}
Erstes Argument ist die Spalte, nicht die Seite. Zwei Zeilen stehen voll,
dann weichen vier Zeilen um 25 mm aus, danach wieder voll. Der Text muss lang
genug sein, dass alle sieben angegebenen Zeilen wirklich entstehen, sonst ist
der Befund nicht aussagekraeftig. Noch ein Satz dafuer.

\dsaabschnitt{Umfluss entlang einer Silhouette}

\dsaUmflussKontur{r}{0}{11.6mm,12.1mm,13.5mm,13.4mm,11.9mm,16.4mm,16mm,17.9mm}
Der Text umfliesst nicht den Bildrahmen, sondern den Freistellpfad. Deshalb
ein Einzug je Zeile. Die acht Werte sind die ersten acht einer gemessenen
Kontur. Der Absatz muss lang genug sein, dass alle acht Zeilen entstehen.

\dsaabschnitt{Tabelle}

\dsaTabellenkopf{Qualitaet und Preis einer Herberge}
\begin{dsaTabelle}{ll}
\dsaTabellenrubrik \dsaRubrikschrift Qualitaet & \dsaRubrikschrift Preis\\
\textbf{1} & jaemmerliche Bruchbude, 50 Prozent\\
\textbf{3} & einfache Herberge, Normalpreis\\
\textbf{6} & luxurioese Unterkunft, 400 Prozent\\
\dsaTabellenrubrik \dsaRubrikschrift Anmerkung
  & \dsaAnmerkungsschrift Preise je Nacht und Person\\
\end{dsaTabelle}

\dsaProbe{Sinnesschaerfe (Suchen), erschwert um 2}
\begin{dsaliste}
\item QS 1: die erste Stufe
\item QS 2: die zweite
\item QS 3 und mehr: die dritte
\end{dsaliste}

\clearpage

%%% ==================================================================
\dsakapitel{Seitentypen}
%%% ==================================================================

Neun Typen, aber nur die zeigen sich in einem Musterbogen, die keine eigene
Seite brauchen.

\begin{dsaSeiteOhneZahl}
Seite ohne Seitenzahl. Hintergrund ja, Seitenzahl nein.
\end{dsaSeiteOhneZahl}

\dsaGanzseite{grafiken/aventurienkarte}

Ganzseitige Grafik im Grafikbereich, 184 mal 265 mm — der zweite Rahmen des
Baukastens.

% Querformat zuletzt, weil \newgeometry das Raster fuer die Folgeseiten neu
% setzt und danach geprueft werden muss.
\dsaQuerAnfang
Querformat, etwa fuer einen Meisterschirm. Nicht im Baukasten, aber ein
haeufiger Bedarf.
\dsaQuerEnde

Text nach der Querseite. Mit \verb|rasterzeigen| nachsehen, ob die
Grundlinien wieder stimmen.

%%% ==================================================================
%%% Rueckseite mit Regionalkarte
%%% ==================================================================
% Die letzte Seite eines Hefts, aus dem Rueckseiten-Karten-Paket. Die Karte
% steht in Sepia, nur die genannte Region in Farbe, mit weichem
% Schlagschatten -- so wie die offiziellen Hefte es zeigen. Hier zum
% Vergleich neben \dsaUmschlagHinten; ein wirkliches Heft nimmt eins von
% beiden.
%
% Zum Pruefen: die Region muss farbig sein, der Rest sepia, und der
% Zierrahmen unveraendert bunt. Wer eine andere Region will, setzt
% ruecken-thorwal, ruecken-bornland und so weiter ein -- 28 stehen bereit.
\dsaRueckseite{ruecken-mittelreich}{Ein Beispiel}{von [Name]}{%
  Hier steht der Klappentext. Absaetze werden durch Leerzeilen getrennt und
  im Blocksatz gesetzt, 81,7 mm breit.

  Der zweite Absatz zeigt den Abstand von 2,83 mm, den das Vorbild
  zwischen zwei Absaetzen laesst.
}{%
  \dsaRueckKopf{Ein DSA-Gruppenabenteuer\\fuer 3 bis 5 Helden}
  \dsaRueckFeld{Genre}{Beispiel}
  \dsaRueckFeld{Ort}{irgendwo}
  \dsaRueckStrich
  \dsaAnforderungen{1}{4}{2}{1}
}

%%% ==================================================================
%%% Umschlag hinten
%%% ==================================================================

\dsaUmschlagHinten{grafiken/umschlag-hinten}{Ein Beispiel}{%
  Hier steht der Klappentext.

  \dsaZusammenfassung{ein Beispiel in wenigen Worten}{Beispiel}{keine}%
    {irgendwo}{jederzeit}

  \dsaAnforderungen{1}{4}{2}{1}
}

\end{document}
